% ============ 第2章 §2.12 Riesz兄弟定理(原书 pp.154-155) ============
\section{Riesz 兄弟定理}\label{sec:2.12}
\index{Riesz 兄弟定理}


本节考察 Cauchy--Stieltjes 积分何时属于 $H^{1}$.这与测度论意义上令人惊奇的 Riesz 兄弟(F. and M. Riesz)定理密切相关.在 2.11 节中我们注意到,属于 $\bigcap_{p<1} H^{p}$ 的函数未必有 Poisson--Stieltjes 积分表示,而定理 2.12.1 表明对于 $H^{1}$ 这是可能的.本节中 $\hat{\mu}$ 表示 $\mu$ 的 Fourier 系数.

\begin{lemma}\label{lem:2.12.1}
设 $\mu \in M[-\pi,\pi]$ 为复测度且 $\hat{\mu}(n) = 0$($n < 0$),则 $\mu$ 的 Poisson--Stieltjes 积分属于 $H^{1}$,而且实际上是 Poisson 积分.
\end{lemma}

\begin{proof}
设 $F(z)$ 为 $\mu$ 的 Cauchy--Stieltjes 积分,则当 $|z| < 1$ 时
\begin{equation*}
\frac{1}{1 - ze^{-it}} = \sum_{n=0}^{\infty} e^{-int}z^{n}
\end{equation*}
且
\begin{equation*}
\frac{e^{it}}{e^{it} - \frac{1}{\bar{z}}} = -\sum_{n=0}^{\infty} e^{i(n+1)t}\bar{z}^{n+1},
\end{equation*}
所以 $F(z) = \sum_{n=0}^{\infty}\hat{\mu}(n)z^{n}$,并且
\begin{equation*}
F\left(\frac{1}{\bar{z}}\right) = -\sum_{n=0}^{\infty}\hat{\mu}(-(n+1))\left(\frac{1}{\bar{z}}\right)^{n+1} \equiv 0.
\end{equation*}
于是当 $|z| < 1$ 时成立
\begin{equation*}
F(z) = F(z) - F\left(\frac{1}{\bar{z}}\right) = \frac{1}{2\pi}\int_{-\pi}^{\pi} P(r, \theta - t)\dd \mu(t),
\end{equation*}
由定理~\ref{thm:2.3.1} 的证明知 $F \in H^{1}$.
\end{proof}

\begin{theorem}[Riesz 兄弟]\label{thm:2.12.2}
若 $\mu \in M[-\pi, \pi]$ 为满足 $\mu(\{-\pi\}) = \mu(\{\pi\})$ 的复测度且 $\hat{\mu}(n) = 0$($n < 0$),则 $\mu$ 关于 Lebesgue 测度绝对连续.
\end{theorem}

\begin{proof}
由引理~\ref{lem:2.12.1},$\mu$ 的 Poisson--Stieltjes 积分 $f(z)$ 属于 $H^{1}$.由定理~\ref{thm:2.9.1},$f(z)$ 可由 $f(e^{it})$ 作 Poisson 积分表示.由 Poisson 积分表示的唯一性(命题~\ref{prop:2.2.3}),$\dd \mu(t) = f(e^{it})\dd t$.
\end{proof}

\begin{proposition}\label{prop:2.12.3}
$f \in H^{1}$ 可由 $f$ 的边界值 $f(e^{it})$ 作 Cauchy 积分表示.
\end{proposition}

\begin{proof}
设 $F$ 为 $f(e^{it})$ 的 Cauchy 积分.令 $\dd \mu = f(e^{it})\dd t$,则由定理~\ref{thm:2.10.5} 有 $\hat{\mu}(n) = 0$($n < 0$).由引理~\ref{lem:2.12.1},$F \in H^{1}$ 且 $F$ 是 Poisson 积分,并由定理~\ref{thm:2.9.1} 等于 $f$.即
\begin{align*}
F(z) & = \frac{1}{2\pi}\int_{-\pi}^{\pi} \frac{e^{it}}{e^{it} - z} f(e^{it})\dd t \\
& = \frac{1}{2\pi}\int_{-\pi}^{\pi} P(r, \theta - t) f(e^{it})\dd t = f(z).
\end{align*}
\end{proof}
